% --- CÁC GÓI CẦN THIẾT BẮT BUỘC (CORE PACKAGES) ---
\usepackage[utf8]{inputenc} % Hỗ trợ gõ tiếng Việt có dấu
\usepackage[T5]{fontenc}    % Định dạng font chữ tiếng Việt
\usepackage{graphicx}       % Chèn hình ảnh
\usepackage{float}          % Cố định vị trí hình/bảng (dùng tham số [H])
\usepackage{tikz}           % Vẽ hình và tạo khung cho trang bìa
\usetikzlibrary{calc}       % Tính toán tọa độ trong TikZ (cần cho khung trang bìa)

% --- GỢI Ý CÁC GÓI MỞ RỘNG (BỎ DẤU % Ở ĐẦU ĐỂ SỬ DỤNG KHI CẦN) ---
% \usepackage[unicode]{hyperref}   % Tạo link nhấp được và bookmark trong file PDF
% \usepackage{amsmath, amssymb}    % Hỗ trợ gõ công thức toán học nâng cao
% \usepackage{geometry}            % Tùy chỉnh lề trang giấy (VD: margin=2.5cm)
% \usepackage{fancyhdr}            % Tạo và tùy chỉnh Header/Footer
% \usepackage{caption}             % Tùy chỉnh chú thích cho hình ảnh và bảng biểu
% \usepackage{subcaption}          % Hỗ trợ chèn các hình ảnh con (subfigures)
% \usepackage{tabularx}            % Tạo bảng tự động điều chỉnh độ rộng cột
% \usepackage{multirow}            % Cho phép gộp nhiều dòng trong bảng
% \usepackage[table]{xcolor}       % Hỗ trợ tô màu chữ và màu nền cho ô bảng
% \usepackage{enumitem}            % Tùy chỉnh ký hiệu và khoảng cách danh sách
% \usepackage{titlesec}            % Tùy chỉnh kiểu dáng tiêu đề chương, mục
% \usepackage{tocloft}             % Tùy chỉnh định dạng mục lục
% \usepackage{indentfirst}         % Tự động thụt lùi dòng cho đoạn văn đầu tiên
% \usepackage{pdflscape}           % Hỗ trợ xoay ngang một số trang giấy
\usepackage[acronym]{glossaries} % Hỗ trợ tạo danh mục từ viết tắt tự động
\makeglossaries
% \usepackage{pgf-umlsd}           % Thư viện hỗ trợ vẽ Sequence Diagram
% \usepackage{tikz-uml}            % Thư viện hỗ trợ vẽ Class/Use-case Diagram
